\chapter{Introduction to SALSA}
\label{chap:Introduction}

Welcome to SALSA!  This short chapter introduces SALSA and the problem domain it is intended to address.

What is SALSA?  Well, it could be a delicious spicy tomato-based sauce, or it may refer to an energetic Afro-Cuban inspired type of dance, but in this case SALSA stands for the \underline{S}urveyor's \underline{A}pplications for \underline{L}east \underline{S}quares \underline{A}djustment.  (Please accept our apologies if that's not what you were hoping for in this user guide.)  A ``survey,'' in very general terms, is a set of spatial measurements that are conducted in order to estimate the relative positions of some number of physical locations, or points.  These spatial measurements may include distances, vertical and horizontal angles, height differences, and 3-D vectors as can be obtained using Global Navigation Satellite Systems (GNSS) data.

The surveyor develops a plan for the collection of these measurements with the objective of obtaining that data which are sufficient to estimate the relative positions of points of interest to a level of precision and a level of confidence as mandated by the particular application.  In general, this plan will include redundancy, for three reasons.  First, repeated measurements of the same quantity (e.g., a distance between two points) will improve the precision of one's estimate of that quantity, per the central limit theorem.  Second, redundant measurements equip us to gauge the self-consistency of the measurements and therefore to make assertions about the uncertainties in our final solution.  Third, redundancy in our measurement design allows us to identify highly discrepant measurements, i.e., blunders, so that we can correct the blunder or exclude it from our solution.

The fact that a survey will include redundancy, and the fact that all measurements include some error, means that our measurements will not agree exactly.  We need a strategy to arrive at a solution (estimates of the locations of our points) that is a ``best fit'' to our measurement data.  This raises two important questions: how are dissimilar measurement types combined, and what constitutes a ``best fit?''

The short list of measurement types mentioned above vary in dimensionality, e.g., distances which are 1-dimensional vs. GNSS vectors which are 3-dimensional, as well as in their native frames of reference, e.g., vertical or zenith angles which are measured relative to the local gravity vector vs. GNSS vectors which are computed and expressed in an Earth-centered Earth-fixed (ECEF) Cartesian reference frame.  We must decide on a common reference frame and then work through the math to express these measurements in this common frame -- a process of rotation and linearization.  Once we build a linear set of equations describing our measurement network, we develop a solution using the principles of least squares.  That is, we minimize the sums of the squares of our (weighted) measurement residuals.  The optimality of the least squares solution is well known, dating to work by Carl Gauss in the early 1800s.  Hence, linearized least squares is in fact the strategy to achieve that ``best fit'' in our complicated nonlinear survey problem.

While the fundamental principles of least squares have not changed in over 200 years, our technology sure has.  Even one generation ago, desktop computers were not powerful enough to compute 3-D least squares solutions for large networks (i.e., hundreds of points).  Instead, adjustments were often reduced to 2-D horizontal solutions and separate 1-D vertical solutions which were then combined in some fashion.  The advent of GPS (yielding inherently 3-D information) and advances in computing power have changed the landscape such that today any modern, rigorous, geodetic least squares adjustment software should execute the solution in 3-D.  SALSA does so, using the Earth-centered Earth-fixed (ECEF) Cartesian frame.  The WGS 84 ellipsoid parameters are used for all conversions between geodetic (latitude, longitude, ellipsoid height) and Cartesian (X, Y, Z) expressions of coordinates.  Chapter \ref{chap:UnderstandingSolver} goes into more detail on how the core LSA solver in SALSA, called \textsf{lsasolver}, operates.


Before we leave this short introductory section, we acknowledge there are many useful references on least squares and geodetic surveying, but there is one in particular that we have found very comprehensive and well written: Adjustment Computations by Charles Ghilani (\cite{ghilani}).  We will cite it often, and we commend it to our readers who wish to gain a deeper understanding of the principles of geodetic least squares adjustment.

